\section{Introduction}
\label{sec:intro}

A streaming hash join keeps one side of the join resident as a partitioned
hash table and probes it with batches arriving from the other side. The
textbook cost model counts comparisons, which suggests that the operator's
throughput should be flat in the number of partitions. Measured on real
hardware it is not: throughput falls by more than a factor of two once the
resident side outgrows the last-level cache, and the fall is entirely
attributable to partitions being evicted between the probes that need
them~\cite{balkesen2013}.

The scheduler is free to choose the order in which buffered probe batches
are dispatched. A partition-oblivious scheduler dispatches in arrival
order, which interleaves partitions arbitrarily and evicts each one
shortly before it is needed again. We show that a scheduler which groups
buffered batches by target partition, and dispatches whole groups,
recovers most of that loss.

Our contributions are:

\begin{itemize}
\item A residency model that predicts, from partition sizes and the
  buffered batch mix alone, how much reordering can recover
  (Section~\ref{sec:method}).
\item A scheduler that applies the model online, with a bounded buffer
  and no change to output order.
\item An evaluation on three join workloads showing a
  \(1.7\times\) mean throughput gain (Section~\ref{sec:results}).
\end{itemize}

The rest of this paper is organized around those three pieces.
